% Section 5.2, from lectures/L05/README.md: what you should be able to do, the questions to test
% yourself with, and the next lecture.

\section{Sammanfattning}
\label{c5:sec:summary}

\needspace{6\baselineskip}
\subsection*{Det här ska du kunna nu}
Efter det här kapitlet och Labb~2 ska du kunna:
\begin{itemize}
  \item Koppla upp ett grindnät med 74HC-kretsar på ett kopplingsdäck, med matning, ingångar och
    utgång, enligt en grindtilldelning och en kopplingslista.
  \item Verifiera ett nät mot dess sanningstabell, rad för rad, med förväntat och uppmätt värde.
  \item Bygga NOT, AND och OR av enbart NAND-grindar, och bygga ett nät med enbart NAND.
  \item Mäta en logisk nivå med en multimeter och avgöra om den är 0, 1 eller ogiltig.
  \item Felsöka ett nät systematiskt: matningen först, sedan signalen grind för grind, och läsa ut
    ur sanningstabellen var felet troligen sitter.
\end{itemize}

\needspace{6\baselineskip}
\subsection*{Frågor att testa dig själv med}
\begin{enumerate}
  \item Varför ska matningen mätas först, även när nätet delvis verkar fungera?
  \item En utgång visar 2,4~V. Vad kan det betyda?
  \item Ett nät är fel på exakt de rader där en viss AND-term ska göra utgången till 1. Var mäter
    du först?
  \item Varför kan en 74HC32 sitta där det ska vara en 74HC08 utan att du märker det på alla rader?
  \item Varför är 74HC02 den krets flest kopplar fel, och vad händer om den kopplas som en 74HC00?
  \item \code{X = AB + C} med enbart NAND har en grind mer än med AND och OR. Vad vinner du ändå?
\end{enumerate}

\needspace{6\baselineskip}
\subsection*{Nästa kapitel}
Kapitel~\ref{c6:ch} handlar om sekvensnät, vippor och mikrodatorns uppbyggnad:
\begin{itemize}
  \item Sekvensnät: kretsar som minns, från reläets självhållning till SR-låset.
  \item D-låset och D-vippan, klockan och flanken.
  \item Register, räknare och skiftregister, byggda av vippor.
  \item Mikrodatorns uppbyggnad på blockschemanivå, och hur varje block går tillbaka till grindar
    och vippor.
\end{itemize}
